\scanpage{107}
Or le deuxième est $\xi_0\underset{\substack{\Vert\\\xi_3}}{\bar\xi_0}\,\xi_1\underset{\substack{\Vert\\\xi_2}}{\bar\xi_1}$, or
\[
 \xi_0\bar\xi_0\overset{\unsure{!}}{=}N_{A_1\otimes A_2/A}(\xi_0)
\]
\[
 \xi_1\bar\xi_1=N_{A_1\otimes A_2/A}\left(\underset{\substack{\Vert\\\sigma_1(\xi_0)}}{\xi_1}\right)
 =\overline{N_{A_1\otimes A_2/A}(\xi_0)}
\]
donc \deleted{\unsure{soit}} $N=N_{A_1\otimes A_2/A}(\xi_0)$, on trouve
\[
 N\bar N=N_{A/O}(N)\qquad\text{OK}.
\]
\figure{\begin{tikzpicture}\draw[->] (0,0)--(1.5,0);\end{tikzpicture}}

On va axiomatiser la situation de \deleted{$\unsure{A_1\ A_1,O}$} $\operatorname{Tr}_{A_1\otimes A_2/A}$ et $N_{A_1\otimes A_2/A}$, \unsure{en oubliant la} situation d'un anneau de base $k$, \inserted{et} d'une ``$k$-algèbre $A$ qui ``se comporte comme si c'était une algèbre quadratique sur $k$'' (i.e. loc. libre de rang $2$). Mais on ne suppose ni que $A$ soit loc. libre \deleted{\unsure{sur}}\inserted{ni même plat} / $k$, ni que
\[
 \alpha:k\longrightarrow A\qquad\lambda\mathrel{-}\lambda.1_A
\]
soit injective. Au lieu de cela, on suppose données deux applications
\[
 \begin{array}{lll}
 t_{A/k}=t:A\longrightarrow k&\text{trace relative},\\
 n_{A/k}=n:A\longrightarrow k&\text{norme relative},
 \end{array}
\]
satisfaisant les conditions suivantes

\textcircled{a} $t$ est linéaire, $n$ est quadratique et $t$ se déduit de $n$, comme \inserted{$t(x)=\varphi_n(x,1)$} $\varphi_n$ forme bil. ass. à $n$
donc on a
\[
 \deleted{\unsure{c'est}}\boxed{n(x+\lambda 1_A)=n(x)+\lambda^2\underset{\substack{\Vert\\1\ \text{cf. ci-dessous b)}}}{n(1)}+\lambda t(x)}
\]
\mspage{106}
